% Generated draft: requires mathematical and source review.
\documentclass[11pt]{article}
\usepackage[margin=25mm]{geometry}
\usepackage{fontspec}
\setmainfont{Latin Modern Roman}
\usepackage{amsmath,amssymb,mathtools,unicode-math}
\setmathfont{Latin Modern Math}
\usepackage{enumitem,needspace,xurl,hyperref}
\setlength{\parindent}{0pt}
\setlength{\parskip}{5pt}
\setlength{\emergencystretch}{3em}
\newcommand{\N}{\mathbb{N}}
\newcommand{\Z}{\mathbb{Z}}
\newcommand{\Q}{\mathbb{Q}}
\newcommand{\R}{\mathbb{R}}
\newcommand{\C}{\mathbb{C}}
\newcommand{\F}{\mathbb{F}}
\newcommand{\A}{\mathbb{A}}
\newcommand{\PP}{\mathbb{P}}
\newcommand{\E}{\mathbb{E}}
\newcommand{\eps}{\varepsilon}
\newcommand{\dd}{\,\mathrm{d}}
\begin{document}
\section*{UnsolvedMath ID: 801135}
\textbf{Source problem code:} catalogue entry
\subsection{Definitions and mathematical statement}
For an integer \(n\geq 2\), let \(\mathbb{F}_n\) be the free group on \(n\) generators. Its group von Neumann algebra is
\[
L(\mathbb{F}_n)=\{\lambda_n(g):g\in\mathbb{F}_n\}'',
\]
where \(\lambda_n(g)\delta_h=\delta_{gh}\) on \(\ell^2(\mathbb{F}_n)\), and the double commutant is taken in \(B(\ell^2(\mathbb{F}_n))\). Equip this \(\mathrm{II}_1\) factor with its canonical normalized trace
\[
\tau_n(x)=\langle x\delta_e,\delta_e\rangle,
\]
where \(e\) is the group identity. Interpret elementary equivalence in the standard continuous first-order language of tracial von Neumann algebras: two such structures are elementarily equivalent if every sentence has the same value in both.

The finite-rank formulation of the supplied question is:
\[
\text{For every }n,m\in\mathbb{N}\text{ with }n,m\geq 2\text{ and }n\ne m,
\qquad
(L(\mathbb{F}_n),\tau_n)\equiv(L(\mathbb{F}_m),\tau_m)?
\]
Equivalently, must \(\varphi^{(L(\mathbb{F}_n),\tau_n)}=\varphi^{(L(\mathbb{F}_m),\tau_m)}\) hold for every sentence \(\varphi\) and every such pair of ranks?

The source does not explicitly specify the range of \(n\) or the logical language. The displayed formulation uses finite nonabelian ranks and the standard tracial continuous language. Whether the intended question also includes the countably infinite rank \(n=\infty\) remains unresolved from the supplied record. Rank \(1\) does not give a \(\mathrm{II}_1\) factor and is excluded from this formulation.
\subsection{Short English statement}
Do all finite-rank nonabelian free group factors satisfy exactly the same continuous first-order sentences as tracial von Neumann algebras? The supplied record calls this the \(\mathrm{II}_1\) factor Tarski problem and labels it open. Its intended rank range and logical language require confirmation.
\subsection{Sources}
\begin{itemize}
\item Local UnsolvedMath registry: permanent ID \(801135\), exact source problem code \texttt{catalogue entry}. The supplied record attributes the problem to Elayavalli, with proposed year 2025, and labels its status as open; this status is supplied metadata, not an independently verified assessment.
\item Elayavalli, \emph{50 Open Problems: Ultraproduct \(\mathrm{II}_1\) Factors} (2025), Problem 36, \url{https://arxiv.org/abs/2511.20377}. Supplied context: \(\mathrm{II}_1\) factor Tarski's problem. The record reports access on 2026-10-16 and extraction by \texttt{arxiv\_latex}; these provenance details are supplied and have not been independently verified.
\end{itemize}
\end{document}
